\documentclass[11pt,a4paper]{article}
\usepackage[utf8]{inputenc}
\usepackage[T1]{fontenc}
\usepackage[margin=0.75in]{geometry}
\usepackage{enumitem}
\usepackage{titlesec}
\usepackage{xcolor}
\usepackage{hyperref}

\definecolor{primarycolor}{RGB}{33,37,41}
\definecolor{secondarycolor}{RGB}{0,102,204}
\definecolor{accentcolor}{RGB}{108,117,125}
\hypersetup{colorlinks=true,urlcolor=secondarycolor}
\titleformat{\section}{\large\bfseries\color{primarycolor}}{}{0em}{}
\titlespacing*{\section}{0pt}{1.0em}{0.4em}
\setlist[itemize]{leftmargin=*,itemsep=0.15em,topsep=0.2em}
\pagestyle{empty}

\begin{document}

\begin{center}
{\Huge\bfseries\color{primarycolor} MARTA KOWAL-BRANDT}\\[0.3em]
{\large QA AUTOMATION ENGINEER | DEVOPS \& CI/CD INFRASTRUCTURE}\\[0.3em]
Wroclaw, Poland | marta.kowalbrandt@example.com | +48 555 303 303 | gitlab.com/mkb-example
\end{center}

\section*{PROFESSIONAL SUMMARY}
QA automation engineer with 7 years of experience and a strong DevOps focus in CI/CD, Kubernetes-based test infrastructure and build tooling. Owns and operates GitLab CI pipelines for three product teams and builds and maintains ephemeral test environments deployed via Helm to a shared Kubernetes cluster. Delivers measurable improvements in pipeline performance and CI runner costs, and develops personal infrastructure projects with k3s, GitOps, monitoring and Terraform. Targeting a move into DevOps or platform engineering.

\section*{SKILLS}
\textbf{CI/CD and containers:} GitLab CI, Jenkins, Docker, Docker Compose, Helm\\
\textbf{Infrastructure and cloud:} Kubernetes (namespace-level access and test-environment deployments), k3s, Terraform (personal project), Argo CD, Ansible (basic), AWS (EC2 and S3 basics; personal learning)\\
\textbf{Monitoring and data:} Prometheus, Grafana, Loki, PostgreSQL, SQL\\
\textbf{Scripting and systems:} Python, Bash, Linux, Git\\
\textbf{Testing:} pytest, Playwright, Selenium, Postman, Java (test code)

\section*{SELECTED PROJECTS}
\textbf{Merge-Request Ephemeral Environments | Odrava Health Software | 2023-2024}
\begin{itemize}
\item Built short-lived environments for each merge request, with the application, PostgreSQL and seeded data deployed via Helm to a shared Kubernetes cluster from GitLab CI.
\item Configured automatic cleanup after merge or 48 hours of inactivity; QA and product teams use the environments for feature review and e2e testing.
\item Technologies: GitLab CI, Helm, Kubernetes, Docker, Bash, PostgreSQL.
\end{itemize}

\textbf{CI Pipeline Speed-up | Odrava Health Software | 2022}
\begin{itemize}
\item Profiled pipelines for three teams, parallelised test jobs, added dependency caching and split e2e tests into a separate stage.
\item Reduced average pipeline time from 38 to 14 minutes and monthly CI runner cost by approximately 30\%.
\item Technologies: GitLab CI, Docker, pytest-xdist.
\end{itemize}

\textbf{Homelab k3s Cluster with GitOps | Personal | 2024-2025}
\begin{itemize}
\item Built a three-node Raspberry Pi 5 k3s cluster, with applications deployed from a Git repository using Argo CD.
\item Set up Prometheus, Grafana and Loki monitoring, Traefik ingress and cert-manager for TLS; provisioned infrastructure for a small Hetzner VPS with Terraform.
\item Technologies: k3s, Argo CD, Prometheus, Grafana, Loki, Terraform, Ansible (basic).
\end{itemize}

\textbf{Flaky Test Dashboard | Odrava Health Software | 2023}
\begin{itemize}
\item Collected CI test results in PostgreSQL and created Grafana dashboards showing flaky tests by team.
\item Flaky-test rate fell from 7\% to below 1\% after teams began addressing the most frequent offenders.
\item Technologies: Python, PostgreSQL, Grafana, GitLab API.
\end{itemize}

\section*{EXPERIENCE}
\textbf{QA Automation Engineer (Test Infrastructure) | Odrava Health Software | Wroclaw | 2021-Present}
\begin{itemize}
\item Owns and operates GitLab CI pipelines for three product teams, approximately 25 developers, covering build, unit, integration and e2e test stages and artifact publishing.
\item Maintains a pytest and Playwright e2e suite of approximately 900 tests; reduced the flaky-test rate from 7\% to below 1\%.
\item Writes Python and Bash tooling for seeding test data into PostgreSQL.
\item Wrote the internal guide to how CI works, used in onboarding, and mentored two QA interns in 2023.
\end{itemize}

\textbf{QA Engineer | Brightmill Software House | Wroclaw | 2019-2021}
\begin{itemize}
\item Tested web and mobile applications for client projects using manual and automated testing.
\item Introduced a Java and Selenium automated regression suite for the largest client project and configured Jenkins jobs for nightly runs.
\end{itemize}

\textbf{Customer Support Specialist (2nd line) | Brightmill Software House | 2018-2019}
\begin{itemize}
\item Investigated customer issues in logs and databases and reproduced bugs for developers.
\end{itemize}

\section*{EDUCATION}
BEng Computer Engineering | Silesian Lowlands University of Technology | 2018

\section*{CERTIFICATION AND PROFESSIONAL DEVELOPMENT}
ISTQB Foundation Level | 2019\\
Preparing for the Certified Kubernetes Administrator (CKA) exam, scheduled for January 2027

\end{document}